\begin{theorem}[A four-unit residual margin]
\label{thm:residual-four-margin}
Every nonempty strongly connected arc-minimal all-deficient oriented graph on
$n$ vertices satisfies
\[
 \sum_x\eta(x)=2M-n\ge4,
 \qquad
 {M\ge\left\lceil\frac n2\right\rceil+2}.
\]
Equivalently,
$|A(D)|\le\lfloor n(n-2)/2\rfloor-2$.
\end{theorem}
\begin{proof}
The established bad-edge witness lemma implies that the set
$Z=\{x:\eta(x)>0\}$ contains a missing pair, and therefore the
total residual is at least two. It also implies that every failure
source of every bad missing edge lies in $Z$.

If the total residual is two, there are exactly two positive vertices,
both of residual one. If it is three and has three positive vertices,
all three likewise have residual one. In either case $|Z|\le3$, and
Corollary~\ref{cor:eta-one-budget} gives
$t(z)\le1\le g(z)+1$ for every $z\in Z$.
The small-failure-source transfer lemma
(Lemma~\ref{lem:small-failure-sources}) therefore gives a Seymour
vertex, a contradiction. This application uses the transfer lemma's
canonical convenient orientations and does not assume that an
arbitrary convenient completion preserves all good missing edges.

It remains to exclude total residual three with
$Z=\{u,v\}$, $\eta(u)=2$, and $\eta(v)=1$. The pair $uv$ is missing.
There is a bad missing edge, since otherwise the good-edge completion
argument would give a Seymour vertex. Its two directional failure
sources must be $u,v$. Consequently
$\Nout(v)\cap U(u)\ne\varnothing$.
Apply Lemma~\ref{lem:eta-two-dichotomy}.

If $t(u)\le2$, both positive vertices satisfy the transfer threshold
$t(z)\le g(z)+1$, so
Lemma~\ref{lem:small-failure-sources} again gives a contradiction.
Otherwise that dichotomy supplies convenience of $v\to u$.
Choose the reference linear extension of the protected relation with
$v$ before $u$. This is possible because missing partners are
incomparable in that relation; a protected chain would, by
transitivity, force an arc between them. The canonical convenient
orientation of $uv$ is then $v\to u$: it is either the only convenient
direction or the direction selected by the reference order.
In the transfer lemma's internal missing-edge digraph on $Z$, the only
vertex with positive out-degree is now $v$, and
$t(v)\le1\le g(v)+1$. The vertex $u$ is a zero-out-degree vertex;
the transfer construction directs bad missing edges to avoid its
failures and requires no bound on $t(u)$. Thus the lemma applies once
more, a contradiction.

The total residual is an integer, so it is at least four. Finally
$2M-n$ has the parity of $n$: for odd $n$ it is at least five.
Both parity cases give the displayed missing-pair bound.
\end{proof}
